Reproducible and verifiable builds increase trust in distributed software artifacts by enabling independent parties to detect artifacts produced by compromised build or release pipelines. 
However, artifact verification requires more than deterministic builds: a verifier must also recover the source state, build environment, dependencies, and build instructions that produced the artifact. 
Decentralized-build ecosystems make this difficult because artifacts are produced through heterogeneous tools, maintainer-controlled workflows, and fragmented metadata. 
As a result, it remains unclear how often artifacts in these ecosystems can be independently verified.

This paper studies artifact verifiability across four popular decentralized-build package ecosystems. 
We define an independent verifier model that relies only on registry-derivable metadata and an artifact comparison model with tiered equivalence levels. 
We implement these models in an Artifact Verification Pipeline and use it to measure artifact verifiability across the target ecosystems. 
Our results show that, beyond build determinism, verifiability is limited by missing source and build metadata, implicit release transformations, and unconventional build practices. 
Provenance attestations and embedded VCS metadata improve verification, but they do not provide complete rebuild specifications. 
These findings identify concrete metadata gaps and ecosystem-level changes needed to make artifact verification practical at package-registry scale.
